\section*{Contribution Declaration}

\begin{table}[h]
  \centering\small
  \setlength{\tabcolsep}{5pt}
  \renewcommand{\arraystretch}{1.2}
  \caption*{\textbf{Individual Contribution Split (\%)}}
  \begin{tabular}{lccc}
    \toprule
    \textbf{Task} & \textbf{Dang} & \textbf{Hieu (pd)} & \textbf{Hieu (pm)} \\
    \midrule
    Data preparation         & 20 & 50 & 30 \\
    Model integration        & 40 & 35 & 25 \\
    Experimentation          & 30 & 40 & 30 \\
    Tactical analysis        & 55 & 15 & 30 \\
    Visualisation            & 25 & 20 & 55 \\
    Manuscript writing       & 35 & 35 & 30 \\
    \midrule
    \textbf{Overall}         & \textbf{34} & \textbf{33} & \textbf{33} \\
    \bottomrule
  \end{tabular}
\end{table}

\section{Conclusion}

We presented \method, an integrated badminton video analytics and data-mining
system linking rally filtering, shuttle tracking, player pose extraction,
court-normalized feature construction, stroke classification, and five
interpretable mining analyses. On unseen ShuttleSet videos~41--44, the clean
evaluation achieves 69.04\% side-aware top-1 accuracy and 87.89\% top-3
accuracy (top-3 measures whether the correct stroke falls within the
model's three highest-scoring classes, a relevant threshold because the
dominant errors involve tactically adjacent strokes---net shot, return
net, and push---that are visually similar and occupy the same forecourt
zone). More importantly, the predicted \texttt{none} rate falls to 6.70\%
and transition-matrix error decreases substantially relative to earlier
baselines, allowing the recovered stroke stream to capture a coherent
attack–defence grammar of net pressure, forced lifts, attacking smashes, and
defensive replies. The five analyses—distributions, transitions, motifs,
movement profiles, and clustering—provide complementary evidence that the
pipeline produces a readable tactical narrative from raw video. The clearest
path forward is to improve far-side player detection and automate the steps
currently supplied by annotations.